\chapter{Canonical Quantization}
\label{cq}
{
    \setlength\epigraphwidth{0.6\textwidth}
\epigraph{ Quantum mechanics is a
different kind of a theory to represent the world.
} {R. P. Feynman}
}
\section{Introduction}
The relation between classical and quantum mechanics is often
described as one of extension: quantum mechanics generalizes
classical mechanics. Actually it is correct only for a
part (though very important one) of classical
 systems ---  systems with phase space $ \mathbb{R}^{2 n} $
 and  Hamiltonian which is quadratic function of momenta,
 and beside them some relatively exotic highly symmetric systems. 
 In that case quantum mechanics
replaces classical mechanics by altering its fundamental
mathematical structure. Classical mechanics reappears only 
as some kind of  singular limit.
On the other hand  crucial components of classical theory 
like constraints, canonical transformations, trajectories,
have no faithful quantum analogues.

Non-relativistic quantum mechanics can be presented axiomatically,
without explicit reference to a classical precursor. 
However all rigorous definitions and theorems give us no way to connect 
the mathematical constructions to the real world. 

Quantum mechanics replaces classical mechanics by altering:
\begin{itemize}
\item The algebra of observables (commutative $\to$ noncommutative)
\item The notion of state (phase space point $\to$ Hilbert space vector)
\item The dynamical bracket (Poisson $\to$ commutator)
\end{itemize}
Logically we should start from some quantum system and 
investigate how and under what conditions classical system emerges.  
"Of course, nature knows the quantum mechanics,
and the classical mechanics is only an approximation; so there is no mystery in
the fact that in classical mechanics there is some shadow of quantum mechanical
laws---which are truly the ones underneath. To reconstruct the original object
from the shadow is not possible in any direct way, but the shadow does help
you to remember what the object looks like" \cite{flp3}. 
This is the physical idea behind quantization which
is the collection of procedures by which a classical mechanical 
system is associated with a non-relativistic quantum theory. 
Naturally, quantization is neither unique nor algorithmic in general 
and no single scheme applies uniformly to all classical systems.
Below  we analyze some quantization schemes, 
discuss no-go theorems limiting the existence of consistent quantization maps.

\section{Bra-Ket Notation}
\label{cq:bk}

Bra--ket notation is the standard symbolic language of quantum
mechanics. While extraordinarily powerful and conceptually elegant,
it compresses multiple layers of mathematical structure into compact
expressions. 
Its success stems from its ability to encode:
\begin{itemize}
\item Linear algebraic structure,
\item Operator theory,
\item Spectral decomposition,
\item Measurement theory,
\item Tensor products and entanglement.
\end{itemize}
However, the notation often suppresses functional-analytic and
geometric subtleties. 

It appears at Paul Dirac's book \cite{dirac-pqm}  "The Principles of Quantum Mechanics"
(first edition 1930). Let's describe pros and cons:
\subsection{ Pros}
\begin{enumerate}
   \item
 Coordinate-free and basis-independent
bra--ket notation expresses quantum states abstractly:
\[|\psi\rangle \in \mathcal{H}, \qquad \langle \phi| \in \mathcal{H}^*\]
This makes the formalism:
\begin{itemize}
\item
  independent of representation (position, momentum, spin, etc.)
\item
  cleanly adaptable to changes of basis
\item
  naturally compatible with symmetry arguments
\end{itemize}
Example:
\[\psi(x) = \langle x|\psi\rangle\]
The same state appears as a function only after choosing the position
basis.
\item
 Clear dual structure
  The notation explicitly encodes:
\begin{itemize}
\item
  kets: vectors \(|\psi\rangle\)
\item
  bras: dual vectors \(\langle\psi|\)
\item
  inner product:
\end{itemize}
\[\langle \phi | \psi \rangle\]
This makes conjugation and adjoints transparent:
\[(|\psi\rangle)^\dagger = \langle \psi|\]
\item
elegant operator representation:
\[A|\psi\rangle\]
outer products generate rank-one operators:
\[|\psi\rangle\langle\phi|\]
projection operators:
\[P = |a\rangle\langle a|\]
spectral decomposition:
\[A = \sum_a a |a\rangle\langle a|\]
or (continuous spectrum)
\[A = \int a\, |a\rangle\langle a|\, da\]
This makes the spectral theorem visually intuitive.
\item
Simplifies tensor products ---
composite systems are naturally written:
\[|\psi\rangle \otimes |\phi\rangle\]
Entangled states appear transparently:
\[|\Psi\rangle = \sum_{ij} c_{ij} |i\rangle \otimes |j\rangle\]
\item
Works naturally with PVM and POVM formalism \footnote{ I'll discuss 
   these things later.     }
projection-valued measure:
\[E(\Delta) = \int_\Delta |x\rangle\langle x| dx\]
POVM elements:
\[F_i = M_i^\dagger M_i\]
Bra--ket notation makes these constructions compact and readable.
\end{enumerate}

\subsection{ Cons}
\begin{enumerate}
   \item
Lack of mathematical rigor (without extra care)---
Dirac freely wrote expressions like:
\[\langle x|x' \rangle = \delta(x-x').\]
But:
\begin{itemize}
\item
  \(|x\rangle \notin L^2(\mathbb{R})\)
\item
  delta functions are distributions, not functions
\item
  continuous ``eigenvectors'' are not true Hilbert space elements.
\end{itemize}
Strictly speaking, this requires the theory of rigged Hilbert spaces:
\[\Phi \subset \mathcal{H} \subset \Phi',\]
Without this, the notation hides subtle functional analysis.
   \item
Can hide domain issues ---
unbounded operators (like momentum or Hamiltonians) require domain care:
\[P = -i\hbar \frac{d}{dx}.\]
In bra--ket notation, one may forget:
\begin{itemize}
\item
  domain restrictions
\item
  self-adjoint extensions
\item
  boundary conditions.
\end{itemize}
These are essential in rigorous analysis.
   \item
Ambiguity in infinite dimensions---
in finite dimensions:
\[\langle \psi| = |\psi\rangle^\dagger\]
in infinite dimensions:
\begin{itemize}
\item
  The identification depends on Riesz representation.
\item
  Not all linear functionals are representable without topology
  assumptions.
\end{itemize}
The notation may conceal these subtleties.
   \item
 Overuse leads to symbolic manipulation without understanding.
Some common problems:
\begin{itemize}
\item
  Treating kets like column vectors everywhere.
\item
  Confusing formal identities with rigorous equalities.
\item
  Ignoring operator domains and closure.
\end{itemize}
It encourages formal calculation over structural analysis.
\end{enumerate}


Dirac's book \cite{dirac-pqm} contains
statements that are mathematically unclear, incomplete, or (by modern standards) incorrect. 
But he did not work in the modern functional-analytic framework that developed later.
For example:
\begin{itemize}
   \item He freely used “eigenvectors” of operators with continuous spectrum.
    \item He manipulated $ \delta $-functions before distribution theory existed.
    \item He treated unbounded operators without domain analysis.
    \item He wrote symbolic identities like \[ \langle x|p \rangle=\exp(i p x/\hbar) \]
as if everything were finite-dimensional.
\end{itemize}
From today’s perspective, this is not rigorous.
From 1930’s perspective, it was revolutionary.

When Dirac wrote his book:
\begin{itemize}
   \item Hilbert space theory was still developing.
\item Spectral theorem was not yet standard physics language.
\item Rigged Hilbert spaces did not exist.
\item Distribution theory did not exist.
\end{itemize}
He was not misusing established mathematics.
He was inventing a new symbolic calculus before the mathematics to justify it existed.
Later, mathematicians 
provided rigorous frameworks that made Dirac’s manipulations meaningful.

However some of statements are genuinely incorrect, for example:
\begin{itemize}
   \item implicit assumption that every observable has a complete set of eigenvectors;
\item overuse of spectral decompositions without domain care;
\item loose treatment of continuous spectra;
\item ambiguities about self-adjoint vs symmetric operators;
\item (most serious!) statement that canonical transformations in classical mechanics 
    may be lifted to unitary transformations in quantum mechanics.
\end{itemize}

In his book \cite{neumann-mfqm}  Mathematische Grundlagen der Quantenmechanik (1932)
von Neumann made everything measure-theoretic and Hilbert-space rigorous and avoided 
$ \delta $-functions entirely.
Though von Neumann’s formulation is mathematically precise,
Dirac’s notation became universal in physics.
Because Dirac’s formalism is vastly more flexible and computationally powerful.

Today we understand that Dirac’s formalism is rigorously captured by:
\begin{itemize}
   \item rigged Hilbert spaces;
\item spectral theorem for unbounded self-adjoint operators;
\item distribution theory
\end{itemize}
and there is no reason not to use it.

\section{Mathematical Structure of Quantum Mechanics}
\label{cq:msqm}

Formally a quantum system is defined by the following definitions and axioms,
describing states,observables,dynamics, and probabilities.
\subsection{States}
    \label{cq:sts}
\begin{definition}
    Pure state is a (non-zero) vector in complex separable Hilbert space $ \mathcal{H} $
\[ |\psi \rangle \in \mathcal{H} .\]
\end{definition}
It is not the most general definition of physical state. Physical state may belong to rigged 
Hilbert space i.e. it may be a limit of sequence of vectors. 
Later we shall discuss more general states (mixed) which are represented by tensor product of
pure states. Meanwhile we shall omit "pure".

\begin{definition}
    The quantity
    \begin{equation}
       \label {cq:pa}
        B(\psi,\phi)=\frac{\langle{\psi},{\phi}\rangle}{\|{\psi}\| \|{\phi}\|}  
    \end{equation} 
    is called probability amplitude.
\end{definition}  
\begin{axiom}[Born rule]
     The probability of the transition between states $ \phi $ and $ \psi $ is
     \begin{equation}
        \label {cq:br}
         P(\psi,\phi)=|B(\psi,\phi)|^2
     \end{equation}
\end{axiom}

\subsection{Observables}
    \label{cq:obs}
\begin{axiom}[Observables]
Observables are self-adjoint operators $ \hat A= \hat A^\dagger$ on $\mathcal H$.
\end{axiom}
Strange thing about quantum mechanics is that observables act on states.
For any observable $ \hat A $ and any (pure) state $ \psi $ exits (pure)
state $ \phi = \hat A \psi $ provided that
\begin{itemize}
    \item $ \psi \in \text{Dom}(A) $
    \item $ \phi \neq 0 $
    \item Super-selection \footnote{Will be discussed later} rules are not violated
\end{itemize}
\begin{remark}
    $  \hat A \psi$ is not a result of measurement of observable  $ \hat A $ in
    state $ \psi $.
\end{remark}
\begin{definition}
    The quantity $ \langle \phi,\hat A \psi \rangle $ is called matrix element.
\end{definition}  
Matrix elements are actually the base of all quantum theory. 
\begin{axiom}[Expectation value]
    Expectation value of an observable in a state 
    \begin{equation}
       \label {cq:exp1}
    \langle \hat A \rangle_\psi = \frac{\langle \psi, \hat A \psi \rangle} {\langle \psi,\psi \rangle}
    \end{equation}
    comprises all available information.
\end{axiom}

\subsection{Dynamics}
    \label{cq:dyn}
    There is an observable $ \hat H$ (Hamiltonian),which generates time evolution
    (unitary group):
    \begin{equation}
       \label {cq:evol}
 \hat U(t) = e^{-i \hat Ht/\hbar}.
    \end{equation}

    Evolution may be seen as action on observables (Heisenberg picture)
\begin{align}
   \hat A(t)&= \hat U(t) \hat A \hat U^{-1}(t), \label {cq:heis1}\\
    \frac{d \hat A}{d t} &= \frac{i}{\hbar}[\hat H,\hat A].\label{cq:heis2}
\end{align}
    Or as action on states (Schr\"odinger picture)
\begin{align}
    |\psi(t)\rangle&= \hat U^{-1}(t) |\psi \rangle, \label{cq:schr1} \\
    \frac{d |\psi \rangle}{d t} &=-\frac{i}{\hbar} \hat H |\psi \rangle. \label{cq:schr2}
\end{align}
   
\begin{remark}
    The structure described above is way too abstract. All infinite
    dimensional separable Hilbert spaces are unitary equivalent to
    $ l^2 $ ---the space of square-summable sequences.
    All commuting  observables can be represented as functions of
    one observable. Such representations certainly have no physical
    sense. So we should use more specific constructions, respecting
    geometry and all that.
\end{remark}

\section{Linear QM vs Non-Linear CM}
\label{cq:lqm-ncm}
Non-linearity in finite dimensional space (phase space of classical mechanics)
can be traded for linearity in infinite dimensions (state space of quantum mechanics).
This slogan has several precise mathematical meanings. It is not mystical 
--- it reflects the fact that many nonlinear evolutions become linear 
when lifted to a suitable function space.

This principle was made systematic by
Bernard Koopman.
Given a nonlinear dynamical system with flow $ \Phi_t $
, define: 
\[ (U_t \phi)(x)= \phi(\Phi_t(x)) \]
Even if $ \Phi_t $
is nonlinear, the operator $ U_t $
  is linear on the function space.
The generator is the linear differential operator:
\[L=F(x) \nabla\]. 
For example instead of evolving points \[ \dot x=F(x) \]
evolve probability densities \[ \partial \rho+\nabla (F \rho)=0 \]
This is the classical Liouville picture.

Actually  nothing becomes simpler or harder --- the difficulty just moves.

\section{Rays or Vectors}
\label{cq:rorv}

Expectation value of an observable in a state \eqref{cq:exp1} does not
depend on length $ \|{\psi}\|  $ and on phase of a vector $ |{\psi}\rangle $.
So in many books on quantum mechanics state is defined not as a vector but
as a ray in Hilbert space which is  equivalence class $ [\psi] $ of 
non-zero vectors with equivalence relation:
\begin{equation}
   \label {cq:eqvrel}
|{\psi}\rangle \sim |{\phi}\rangle \quad 
    \text{iff} \quad |{\psi}\rangle = 
    \lambda |{\phi}\rangle \quad \lambda \in \mathbb{C}, \  \lambda \neq 0
\end{equation}
So states are seen as points in  projective Hilbert space:$\mathcal{P}(\mathcal{H})$.
This approach allows very nice geometric picture,
(see "Geometric Quantum Mechanics below), but since $\mathcal{P}(\mathcal{H})$
is not a vector space, it destroys (or at least severely complicates) the
basic (linear) structures of quantum mechanics:
\begin{itemize}
\item superposition
\item Schr\"odinger dynamics
\item entanglement.
\end{itemize}

Quantum mechanics has inherently a gauge symmetry. In electrodynamics
fields are physical quantities they and only they appear in dynamic 
equations but potentials supply working tools \footnote{ In quantum 
mechanics electromagnetic potential IS and observable}. Similarly
in quantum mechanics expectations and projectors depend only on
rays, but to compute that rays we essentially need vectors.

\section{Dirac Quantization}
\label{cq:dirac-quantization}
However  not all Dirac's ideas prove to be right. He sees
quantization as a map of classical Hamiltonian system into
quantum one.

A classical non-relativistic system is described by a symplectic manifold $(M,\omega)$, 
typically $M=T^*Q$ with canonical symplectic form
\begin{equation}
\omega = dq^i \wedge dp_i.
\end{equation}
Observables are smooth functions $f \in C^\infty(M)$ equipped with the Poisson bracket
\begin{equation}
{f,g} = \omega(X_f,X_g),
\end{equation}
where $X_f$ is the Hamiltonian vector field defined by $\iota_{X_f}\omega = df$. 
The resulting Poisson algebra encodes both kinematics and dynamics.

Time evolution is generated by a Hamiltonian function $H$ according to
\begin{equation}
\frac{df}{dt} = \{f,H\}.
\end{equation}
This structure motivates the canonical correspondence between Poisson brackets and 
commutators in quantum theory.

One often idealizes quantization as a map
\begin{equation}
Q : C^\infty(M) \to \mathcal{A}(\mathcal{H}),
\end{equation}
from classical observables to operators on a Hilbert space $\mathcal{H}$, satisfying 
linearity, reality, normalization, and the Poisson--commutator correspondence
\begin{equation}
    \label{cq:pcc}
Q(\{f,g\}) = \frac{1}{i\hbar}[Q(f),Q(g)].
\end{equation}

However the Groenewold--Van Hove theorem shows that no such map exists on the full 
Poisson algebra of observables for generic phase spaces. 
Even for $M=\mathbb{R}^{2n}$, the correspondence can be consistently imposed only on 
restricted subalgebras \footnote{ In my paper \cite{panfil-q} I studied these subalgebras},
demonstrating that quantization is intrinsically limited.

Suppose that phase space admits global Cartesian coordinates 
and restrict \eqref{cq:pcc} to these variables, 
then quantization assigns Hermitian (self-adjoint) operators
$\hat q_i, \hat p_j$ on Hilbert space $ \mathcal{H} $
satisfying canonical commutation relations
\begin{equation}
    \label{cq:ccr}
    [\hat q_i, \hat p_j] = i\hbar \delta_{ij}.
\end{equation}
to classical observables $q_i, p_j$ --- coordinates on phase space $ \mathbb{R}^{2 n} $.
All other observables are obtained by substitution of these operators into
classical observables --- functions on phase space.
Generally operator ordering ambiguities arise,  
and additional analytic input is required to ensure self-adjointness.
This prescription is insufficient for general systems, it works only
if system is described by global Cartesian coordinates.
So the standard name "canonical quantization" is misleading, but generally used.

For "usual" Hamiltonian function $ H=p^2/(2 m) +V(q) $ we get:
\begin{equation}
   \label {cq:ham}
    \hat H = \frac{\hat p^2}{2 m}+V(\hat q)
\end{equation}

In particular using \eqref{cq:ham},\eqref{cq:ccr} and \eqref{cq:heis2} we get:
\begin{equation}
   \label {cq:ehrenfest}
   \begin{split}
       \frac{d \hat q_i}{d t} &= \frac{\hat p_i}{m}\\
       \frac{d \hat p_j}{d t} &= - \frac{\partial V}{\partial q_j}(\hat q)
   \end{split}
\end{equation}
which formally coincide with classical equation of motion
--- Ehrenfest's theorem.
\begin{remark}
    Spectrum of all operators $ \hat q_i, \hat p_j$ satisfying
    canonical commutation relations \eqref{cq:ccr} is unbounded.
    This is absolutely reasonable since these operators 
    correspond to \emph{global} Cartesian coordinates.
\end{remark}  
\section{Some Mathematics}
\label{cq:math}
 There is some confusion in terminology, exact  operator domain is 
 rarely important in physical problems, so physicists generally do 
 not pay attention to difference between Hermitian and self-adjoint operators.
\begin{definition}
    Operator $ \hat A^\dag$ is called adjoint to operator $ \hat A $,
    if 
$ \langle \hat A\phi,\psi \rangle = \langle \phi,\hat A^\dag \psi \rangle$.
\end{definition}
Note that domain of $ \hat A^\dag$ includes that of $ \hat A $
$ D( \hat A)\subseteq D( \hat A^\dag)$.
\begin{definition}
    Operator $ \hat A $ is called Hermitian if $ \hat A = \hat A^\dag $ on
    $D( \hat A)$.
\end{definition}
\begin{definition}
    Operator $ \hat A $ is called  self-adjoint if $ \hat A = \hat A^\dag $
    and $D( \hat A) = D( \hat A^\dag)$.
\end{definition}
For bounded operators on Hilbert space $ \mathcal{H} $ 
$D( \hat A) = \mathcal{H} = D( \hat A^\dag)$ so Hermitian
and self-adjoint bounded operators are the same. 
In many cases Hermitian operators used in quantum mechanics 
are not self-adjoint but have self-adjoint extensions.

Let $ (M,\mu) $ be a 
measure space and $ \mathcal{H} $ a separable Hilbert space with basis $ \{\phi_k\} $. 
Then each $ g \in L^2(M,d \mu;\mathcal{H}) $
is a limit of finite linear combinations of 
vectors of the form $ f_k(x) \phi_k,\ f_k(x) \in L^2(M,d \mu) $ .
We now define 
\[ 
U: \Sigma_{k=1}^N f_k(x)\otimes \phi_k \to
\Sigma_{k=1}^N f_k(x) \phi_k
\]
Then $ U $ is a well-defined map from a dense set in $ L^2(M,d \mu) \otimes \mathcal{H} $
onto a dense set in $ L^2(M,d \mu;\mathcal{H}) $
which preserves norms, so $ U $ U is  
the natural isomorphism between 
$ L^2(M,d \mu) \otimes \mathcal{H} $ and
$ L^2(M,d \mu;\mathcal{H}) $.

\begin{theorem}
    \label{cq:math-t}
    Let $ L^2(M_1,d \mu_1)$ and $ L^2(M_2,d \mu_2)$
 be  separable Hilbert spaces. Then
    \begin{enumerate}
       \item there is a unique isomorphism from
$ L^2(M_1,d \mu_1) \otimes  L^2(M_2,d \mu_2)$ to
$ L^2(M_1\times M_2,d \mu_1\otimes d \mu_2)$
so that $ f\otimes g \to f g $;
\item
 if $ \mathcal{H} $ is a separable Hilbert space, then there is a unique isomorphism from
$ L^2(M,d \mu) \otimes \mathcal{H} $ to
$ L^2(M,d \mu;\mathcal{H}) $
so that $ f\otimes \phi \to f \phi $;
\item
 There is a unique isomorphism from
$ L^2(M_1\times M_2,d \mu_1\otimes d \mu_2)$ to
$ L^2(M_1,d \mu_1;L^2(M_2,d \mu_2)$ 
such that $f(x, y)$ is taken into the function 
$ x \to f(x,\cdot) $. 
    \end{enumerate}
\end{theorem}

\section{Spin}
\label{cq:spin}

 Spin is a purely quantum phenomenon with discrete states. 
 For a particle with spin quantum number $ s $,(where $s$ can be 0,1/2,1,3/2,\ldots)
 the spin Hilbert space is a 
 finite-dimensional complex vector space of dimension $2 s+1$, denoted 
 as $\mathbb{C}^{2 s+1}$.
The total Hilbert space for one particle, 
with spin $s$ is the tensor product of its spatial Hilbert space and its spin Hilbert space:
\begin{equation}
   \label {cq:sps}
    \mathcal{H}_1=L^2(\mathbb{R}^3)\otimes \mathbb{C}^{2 s+1}
\end{equation}
Here spatial part $L^2(\mathbb{R}^3)$ describes the particle's position or momentum. 
According to Theorem \ref{cq:math-t} $\mathcal{H}_1$ is naturally isomorphic to
$L^2(\mathbb{R}^3, dx; \mathbb{C}^{2 s+1})$ --- the Hilbert space of $ 2 s+1 $--component
wave functions.

In particular the Hilbert space in the quantum-mechanical description 
of a single electron $L^2(\mathbb{R}^3, dx; \mathbb{C}^2)$, that is, the 
set of pairs $\{\psi_1(x),\psi_2(x)\}$ of 
square-integrable functions ($dx$ is Lebesgue 
measure). By theorem \ref{cq:math-t}, 
$L^2(\mathbb{R}^3, dx; \mathbb{C}^2)$ is naturally isomorphic to
$ L^2(\mathbb{R}^3)\otimes \mathbb{C}^2 $.

\section{Many particles}
\label{cq-mp}

If we have N distinguishable particles (e.g., an electron, a muon, and a proton),
the total Hilbert space is simply the tensor product of their individual 
single-particle Hilbert spaces.
\begin{equation}
   \label {cq-mp-d}
    \mathcal{H}_N=\bigotimes \mathcal{H}_1^{(i)} \ i=1,2,\ldots,N
\end{equation}
A general state in this space is a linear combination of basis states of the form
$ |{\phi_1}\rangle \otimes |{\phi_2}\rangle \otimes \ldots \otimes |{\phi_N}\rangle $
$ |{\phi_i}\rangle $ contains both the spatial and spin state of the i-th particle.

In quantum mechanics, particles of the same species (e.g., N electrons) 
are fundamental indistinguishable.
Because of this Nature restricts the actual physical states to specific sub-spaces of 
$\mathcal{H}_N$  based on the Spin-Statistics Theorem.
 Boson (integer spin) states 
 must be completely symmetric under the exchange of any two particles.
If you swap the spatial and spin coordinates of particle i and particle j,
the overall quantum state remains exactly the same.
 Fermion (half-integer spin) states
 must be completely antisymmetric under the exchange of any two particles.
If you swap the spatial and spin coordinates of particle i and particle j, 
the overall quantum state picks up a minus sign.
This antisymmetry is the mathematical origin of the Pauli Exclusion Principle:
two fermions can not occupy the exact same spatial and spin state.
